% Source: https://artofproblemsolving.com/community/c3h2392420_more_geometric_probability
% Stable identity: aops-topic-2392420-q01
Put the central square $C$ at $[0,1]\times[0,1]$, with arms $N=[0,1]\times[1,2]$, $S=[0,1]\times[-1,0]$, $E=[1,2]\times[0,1]$ and $W=[-1,0]\times[0,1]$. The cross has area $5$, so each point lies in each of the five squares with probability $\dfrac15$, independently. Boundary cases (a point on an edge, a segment through a corner) have probability $0$ and are ignored. A segment between two points of the cross crosses the perimeter exactly when it leaves the cross.

\textbf{Pairs that never leave the cross.} If both points lie in the same square, or one lies in $C$ and the other in an arm, or they lie in opposite arms, then both points lie in one rectangle ($C$ with one arm is a $1\times2$ rectangle; $N$, $C$, $S$ together form the $1\times3$ rectangle $[0,1]\times[-1,2]$, and likewise $W$, $C$, $E$). A rectangle is convex, so the segment stays inside it.

\textbf{Adjacent arms.} The remaining case is two adjacent arms, for example $P\in N$ and $Q\in E$. Write
\[P=(1-u,\,1+s),\qquad Q=(1+t,\,1-v),\]
where $u,s,t,v$ are independent and uniform on $[0,1]$: $u$ and $t$ are the horizontal distances of $P$ and $Q$ from the line $x=1$, and $s$ and $v$ their vertical distances from the line $y=1$.

\begin{asy}
usepackage("amsmath");
size(5.5cm);
defaultpen(fontsize(10pt));
path sq(pair p){return shift(p)*((0,0)--(1,0)--(1,1)--(0,1)--cycle);}
fill(sq((0,0)),rgb(0.88,0.92,1)); fill(sq((0,1)),rgb(0.88,0.92,1)); fill(sq((1,0)),rgb(0.88,0.92,1));
fill(sq((0,-1)),rgb(0.95,0.95,0.95)); fill(sq((-1,0)),rgb(0.95,0.95,0.95));
draw((0,-1)--(1,-1)--(1,0)--(2,0)--(2,1)--(1,1)--(1,2)--(0,2)--(0,1)--(-1,1)--(-1,0)--(0,0)--cycle,linewidth(1));
draw((0,0)--(1,0)--(1,1)--(0,1)--cycle,gray);
draw((1,1)--(2,1)--(2,2)--(1,2),dotted);
label("missing",(1.5,1.6),fontsize(8pt)); label("corner",(1.5,1.4),fontsize(8pt));
draw((1,-1.25)--(1,2.25),dashed);
label("$x=1$",(1,-1.25),S,fontsize(9pt));
pair Pp=(0.35,1.7), Qp=(1.75,0.55);
draw(Pp--Qp,red+linewidth(0.9));
dot(Pp); dot(Qp);
label("$P$",Pp,NW); label("$Q$",Qp,SE);
draw(Pp--(1,1.7),dotted); label("$u$",(0.675,1.7),N,fontsize(9pt));
draw(Pp--(0.35,1),dotted); label("$s$",(0.35,1.35),W,fontsize(9pt));
draw(Qp--(1,0.55),dotted); label("$t$",(1.375,0.55),S,fontsize(9pt));
draw(Qp--(1.75,1),dotted); label("$v$",(1.75,0.78),E,fontsize(9pt));
label("$N$",(0.15,1.12),fontsize(8pt)); label("$E$",(1.12,0.12),fontsize(8pt)); label("$C$",(0.12,0.12),fontsize(8pt));
\end{asy}

The segment $PQ$ crosses the line $x=1$ at the point dividing $PQ$ in the ratio $u:t$, namely $\dfrac{tP+uQ}{u+t}$, whose height is
\[y^*=\dfrac{t(1+s)+u(1-v)}{u+t}=1+\dfrac{st-uv}{u+t}.\]
\begin{itemize}
\item If $y^*\le1$: the part of the segment with $x\le1$ has heights between $1-v\ge0$ and $1+s\le2$, so it lies in the column $C\cup N$; the part with $x\ge1$ has heights between $y^*$ and $1-v$, both in $[0,1]$, so it lies in $E$. The segment stays inside the cross.
\item If $y^*>1$: just to the right of the line $x=1$ the segment is still above the line $y=1$, so it passes through the missing corner square $(1,2)\times(1,2)$ and leaves the cross.
\end{itemize}
So, given $P\in N$ and $Q\in E$, the segment leaves the cross exactly when $st>uv$. The products $st$ and $uv$ are independent and have the same distribution, so by symmetry $P(st>uv)=P(uv>st)$; and $P(st=uv)=0$. Hence this conditional probability is $\dfrac12$.

By symmetry the same holds for each of the $4$ pairs of adjacent arms $\{N,E\},\{E,S\},\{S,W\},\{W,N\}$, with either point in either arm. The probability that the two points lie in a given pair of adjacent arms (in either order) is $2\cdot\dfrac15\cdot\dfrac15=\dfrac{2}{25}$. Therefore the required probability is
\[4\cdot\dfrac{2}{25}\cdot\dfrac12=\boxed{\dfrac{4}{25}}.\]
